% Source: https://artofproblemsolving.com/community/c3h2740727_quadratic_absolute_value_proof
% Stable identity: aops-topic-2740727-q01
The equation means $x^2-bx+(c-d)=0$ or $x^2-bx+(c+d)=0$, with discriminants $\Delta_1=b^2-4c+4d$ and $\Delta_2=b^2-4c-4d\le\Delta_1$. When $d>0$, no number solves both, since $x^2-bx+c$ cannot equal both $d$ and $-d$. A quadratic with positive discriminant has two real roots with sum $b$; with zero discriminant it has one root $\dfrac b2$.

So for $d>0$ the sum of the distinct real solutions is:
\begin{itemize}
\item $2b$ if $\Delta_2>0$;
\item $\dfrac{3b}{2}$ if $\Delta_2=0$;
\item $b$ if $\Delta_2<0<\Delta_1$;
\item $\dfrac b2$ if $\Delta_1=0$;
\item $0$ (no solutions) if $\Delta_1<0$.
\end{itemize}
For $d=0$ it is $b$, $\dfrac b2$ or $0$ according as $b^2-4c$ is positive, zero or negative.

The claim is \textbf{false}. For example, $\left|x^2+2x\right|=1$ ($b=-2$, $c=0$, $d=1$) has solutions $-1\pm\sqrt2$ and $-1$, with sum $-3\ne4$. Even with four solutions the sum is $2b$, not $2|b|$: $\left|x^2+5x\right|=1$ has sum $-10$.
